\chapter{GitHubによる研究成果管理}
\label{chap8}

\begin{chapterbox}
\textbf{【この章で学ぶこと】}
\begin{itemize}
\item Git/GitHubの基本操作を習得する
\item AIを活用したWebページ作成とGitHub Pagesでの公開を実践できる
\item 研究ポートフォリオの構築と活用方法を理解する
\end{itemize}
\end{chapterbox}
\vspace{5mm}

\section{なぜ大学院生がGitHubを使うべきか}
\index{GitHub@GitHub}

\subsection{研究コードのバージョン管理}

大学院での研究では、コードを何度も修正する。バージョン管理\index{ばーじょんかんり@バージョン管理}なしでは、以下のような問題が生じる。

\begin{verbatim}
よくある悲劇:
analysis.py
analysis_v2.py
analysis_v2_final.py
analysis_v2_final_revised.py
analysis_v2_final_revised_REAL_FINAL.py
\end{verbatim}

このようなファイル管理では、「どの時点で何を変更したか」「なぜ変更したか」が全く分からなくなる。Gitを使えば、コミットした変更の履歴を追跡でき、いつでも過去の状態に戻れる。

\subsection{共同研究での活用}

研究は一人で行うものではない。GitHubを使えば、次のようなことができる。

\begin{itemize}
\item \textbf{コードの共有}: URLを送るだけで共同研究者がコードにアクセスできる
\item \textbf{同時作業}: 複数人が同じプロジェクトを同時に編集できる
\item \textbf{変更の追跡}: 誰が、いつ、何を変更したかが一目で分かる
\item \textbf{レビュー}: 共同研究者の変更をレビューしてからマージできる
\end{itemize}

\subsection{就職活動でのポートフォリオ効果}

技術系の就職活動において、GitHubのプロフィールは「第二の履歴書」になる。特にIT企業やデータサイエンス関連のポジションでは、GitHubでの活動が評価される。

\begin{itembox}[l]{GitHubプロフィールで評価されるポイント}
\begin{itemize}
\item 定期的なコミット（継続的に開発・研究している証拠）
\item README.md が整備されたリポジトリ（ドキュメンテーション能力）
\item コードの品質（命名規則、コメント、テストの有無）
\item オープンソースプロジェクトへの貢献
\end{itemize}
\end{itembox}

\subsection{オープンサイエンスへの貢献}
\index{おーぷんさいえんす@オープンサイエンス}

近年、研究のオープン化（Open Science）が世界的なトレンドとなっている。論文のデータや分析コードをGitHubで公開することで、以下のメリットがある。

\begin{itemize}
\item \textbf{研究の透明性向上}: 分析手法が第三者に検証可能になる
\item \textbf{再現性の担保}: 他の研究者が同じ分析を再現できる
\item \textbf{引用の増加}: 公開されたコードやデータは他の研究者に利用され、引用が増える傾向がある
\item \textbf{学術コミュニティへの貢献}: 研究ツールやデータセットの共有
\end{itemize}

\section{Gitの基礎}
\index{Git@Git}

\subsection{バージョン管理とは何か}

バージョン管理とは、ファイルの変更履歴を記録し、過去の任意の時点の状態を復元できるようにする仕組みである。

\subsection{リポジトリ、コミット、ブランチの基本概念}

Gitを理解するための3つの核心概念を、ゲームのアナロジーで説明する。

\begin{itembox}[l]{Gitの基本概念}
\index{りぽじとり@リポジトリ}\index{こみっと@コミット}\index{ぶらんち@ブランチ}

\textbf{リポジトリ（Repository）= セーブデータフォルダ}

プロジェクトの全ファイルとその変更履歴を格納する場所である。

\textbf{コミット（Commit）= セーブポイント}

ある時点でのファイルの状態を保存する操作である。各コミットには「何を変更したか」を説明するメッセージを付ける。

\textbf{ブランチ（Branch）= 並行世界}

メインの開発ラインから分岐して、独立に作業できる仕組みである。
\end{itembox}

\begin{table}[H]
\centering
\caption{Git用語とゲームのアナロジー}
\begin{tabular}{lll}
\toprule
Git 用語 & ゲームのアナロジー & 説明 \\
\midrule
リポジトリ & セーブデータフォルダ & プロジェクト全体の格納場所 \\
コミット & セーブポイント & ある時点の状態を記録 \\
ブランチ & 並行ルート & 分岐して独立に作業 \\
マージ & ルートの合流 & 分岐した作業を統合 \\
プッシュ & クラウドに同期 & ローカルの変更をGitHubに送る \\
プル & クラウドから同期 & GitHubの変更をローカルに取得 \\
クローン & セーブデータをコピー & GitHubからプロジェクトをダウンロード \\
\bottomrule
\end{tabular}
\end{table}

\section{GitHubの基本操作}

\subsection{GitHubアカウントの作成}

\begin{enumerate}
\item \texttt{https://github.com} にアクセスする
\item 「Sign up」をクリックし、メールアドレス・ユーザー名・パスワードを入力する
\item メール認証を完了する
\end{enumerate}

ユーザー名は本名またはそれに近い名前が望ましい（就職活動で使うため）。GitHub Education（\texttt{https://education.github.com}）に大学のメールアドレスで学生認証を行うと、GitHub Copilotが無料で使えるなどの特典がある。

\subsection{GitHub Desktop のセットアップ}
\index{GitHub Desktop@GitHub Desktop}

コマンドラインが苦手な人には、GitHub Desktop（GUIアプリ）が推奨される。\texttt{https://desktop.github.com} からダウンロードし、GitHubアカウントでサインインする。

\subsection{リポジトリの作成とクローン}

GitHub上でリポジトリを作成する手順:
\begin{enumerate}
\item GitHub にログイン
\item 右上の「+」→「New repository」をクリック
\item Repository name、Description を入力し、Public または Private を選択
\item 「Add a README file」にチェックし、「Create repository」をクリック
\end{enumerate}

\subsection{コミット→プッシュの基本フロー}

日常的な作業の流れは以下の通りである。

\begin{enumerate}
\item ファイルを編集する（VS Code等で）
\item 変更をステージングする（保存したい変更を選択）
\item コミットメッセージを書く
\item コミットする（ローカルにセーブ）
\item プッシュする（GitHubにアップロード）
\end{enumerate}

\begin{itembox}[l]{良いコミットメッセージの書き方}
\begin{table}[H]
\centering
\begin{tabular}{ll}
\toprule
悪い例 & 良い例 \\
\midrule
「更新」 & 「t検定の結果を棒グラフで可視化」 \\
「修正」 & 「欠損値処理のバグを修正: NaNが残る問題」 \\
「追加」 & 「相関分析のコードとヒートマップを追加」 \\
\bottomrule
\end{tabular}
\end{table}
コミットメッセージは「何を、なぜ変更したか」が分かるように書く。
\end{itembox}

\subsection{VS Code との連携}

VS Code\index{VS Code@VS Code}にはGitの操作機能が組み込まれている。左サイドバーの「ソース管理」アイコンからファイルのステージング、コミット、プッシュが可能である。拡張機能として\textbf{GitLens}（コードの各行の変更履歴表示）や\textbf{Git Graph}（コミット履歴のグラフ表示）の導入を推奨する。

\section{AIでWebページを作成する}
\index{Webページ@Webページ}

\subsection{AIでHTML/CSSコードを生成する}

AIを使えば、プログラミングの知識がなくても自己紹介Webページを作成できる。

\begin{promptbox}
大学院生の自己紹介Webページを作成するHTML/CSSコードを書いてください。

内容:
\begin{itemize}
\item 名前: 田中太郎
\item 所属: ○○大学大学院 情報科学研究科 修士2年
\item 研究テーマ: 「深層学習を用いた自然言語処理の研究」
\item スキル: Python, PyTorch, 統計分析
\end{itemize}

デザイン要件:
\begin{itemize}
\item モダンでシンプルなデザイン
\item レスポンシブ対応（スマホでも見やすい）
\item セクション: ヘッダー、自己紹介、研究概要、業績、スキル、連絡先
\item CSSは同一ファイルに含める（1ファイルで完結）
\end{itemize}
\end{promptbox}

\subsection{VS Codeでの編集とプレビュー}

生成されたHTMLをVS Codeで編集する手順は以下の通りである。

\begin{enumerate}
\item VS Codeでプロジェクトフォルダを開く
\item \texttt{index.html}という名前でファイルを作成
\item AIが生成したコードを貼り付ける
\item 拡張機能「Live Server」をインストール
\item \texttt{index.html}を右クリック→「Open with Live Server」
\item ブラウザでリアルタイムにプレビューが表示される
\end{enumerate}

\subsection{AIでカスタマイズする}

Webページの修正やカスタマイズもAIに依頼できる。デザインの変更やセクションの追加など、自然言語で指示するだけで対応してもらえる。

\section{GitHub Pagesで公開する}
\index{GitHub Pages@GitHub Pages}

\subsection{GitHub Pagesとは}

GitHub Pagesは、GitHubのリポジトリに置いたHTMLファイルを無料でWebサイトとして公開できるサービスである。

\begin{itemize}
\item 無料で利用可能
\item \texttt{https://ユーザー名.github.io/リポジトリ名/} のURLで公開される
\item リポジトリにプッシュするだけで自動的に更新される
\item カスタムドメインも設定可能
\end{itemize}

\subsection{GitHub Pages の設定手順}

\begin{enumerate}
\item GitHubで新しいリポジトリを作成する（Publicに設定）
\item リポジトリに \texttt{index.html} をアップロードする
\item リポジトリの「Settings」→「Pages」をクリック
\item 「Source」でBranch: \texttt{main}、Folder: \texttt{/ (root)} を設定し「Save」
\item 数分後に \texttt{https://ユーザー名.github.io/リポジトリ名/} で公開される
\end{enumerate}

\subsection{更新フロー}

ローカルで\texttt{index.html}を編集し、Live Serverでプレビュー確認した後、GitHub Desktopでコミット・プッシュすると、数分後にWebサイトに反映される。

\section{研究ポートフォリオの構築}
\index{ぽーとふぉりお@ポートフォリオ}

\subsection{ポートフォリオに含めるもの}

\begin{table}[H]
\centering
\begin{tabular}{llc}
\toprule
セクション & 内容 & 重要度 \\
\midrule
自己紹介 & 名前、所属、研究分野、顔写真 & 必須 \\
研究概要 & 研究テーマの説明（一般向け） & 必須 \\
業績一覧 & 論文、学会発表、受賞歴 & 必須 \\
スキル & プログラミング言語、ツール、資格 & 推奨 \\
プロジェクト & 研究プロジェクトの詳細と成果物 & 推奨 \\
連絡先 & メールアドレス、SNS & 必須 \\
CV/履歴書 & PDF形式でダウンロード可能 & 推奨 \\
\bottomrule
\end{tabular}
\end{table}

\subsection{README.md の書き方}

README.md\index{README@README.md}は、リポジトリの「顔」である。GitHubでリポジトリを開いたとき最初に表示されるため、丁寧に書くことが重要である。プロジェクト概要、特徴、インストール方法、使い方、ファイル構成、依存ライブラリ、ライセンス、著者情報を含めるとよい。

\subsection{フォルダ構成のベストプラクティス}

\begin{verbatim}
my-research-project/
├── README.md
├── LICENSE
├── requirements.txt
├── .gitignore
├── src/              # ソースコード
├── data/
│   ├── raw/          # 生データ（変更不可）
│   └── processed/    # 処理済みデータ
├── notebooks/        # Jupyter Notebook
├── figures/          # 生成されたグラフ
├── tests/            # テストコード
└── docs/             # ドキュメント
\end{verbatim}

\begin{itembox}[l]{.gitignore ファイル}
\texttt{.gitignore}は、Gitで追跡すべきでないファイルを指定するファイルである。大きなデータファイル、Pythonのキャッシュ（\texttt{\_\_pycache\_\_/}）、環境変数（\texttt{.env}）、OS固有ファイル（\texttt{.DS\_Store}）などを指定する。
\end{itembox}

\section{共同研究でのGitHub活用}

\subsection{Issue による課題管理}
\index{Issue@Issue}

Issueは、プロジェクトの課題やタスクを管理する機能である。タイトルと説明を記入し、ラベル（bug, enhancement, question等）を付与し、担当者を指定する。研究プロジェクトでは、データ前処理、統計分析、図表作成などのタスクをIssueとして管理できる。

\subsection{Pull Request によるコードレビュー}
\index{Pull Request@Pull Request}

Pull Request（PR）は、ブランチの変更をメインブランチに統合する前に、レビューを行う仕組みである。

\begin{enumerate}
\item 新しいブランチを作成する
\item ブランチ上でコードを書いてコミットする
\item GitHub上でPull Requestを作成する
\item 共同研究者がコードをレビューする
\item レビューが通ったらマージする
\end{enumerate}

\subsection{研究コードの共有とライセンス}
\index{らいせんす@ライセンス}

研究コードを公開する場合、適切なライセンスを選択することが重要である。

\begin{table}[H]
\centering
\begin{tabular}{lll}
\toprule
ライセンス & 特徴 & 推奨場面 \\
\midrule
MIT & 最も自由。商用利用・改変・再配布可能 & 研究コードの一般公開 \\
Apache 2.0 & MIT + 特許条項 & ツール・ライブラリの公開 \\
GPL v3 & 派生物も同じライセンスで公開必須 & 強いオープンソースを望む場合 \\
CC BY 4.0 & クリエイティブ・コモンズ & データセットの公開 \\
\bottomrule
\end{tabular}
\end{table}

研究コードには\textbf{MITライセンス}が最も広く使われている。特別な事情がなければMITが無難な選択肢である。

\section*{演習問題}

\begin{enumerate}
\item[\textbf{演習8-1}] \textbf{GitHubアカウントとリポジトリの作成}\\
GitHubアカウントを作成し、GitHub Educationに登録しなさい。\texttt{my-first-repo}という名前のリポジトリを作成し、GitHub Desktopでクローンした後、テキストファイルを1つ追加してコミット・プッシュしなさい。スクリーンショットを撮影して提出すること。

\item[\textbf{演習8-2}] \textbf{AIでWebページを作成する}\\
AIチャットツールを使って、自己紹介Webページを作成しなさい。自分の情報を含むプロンプトを作成し、AIにHTML/CSSコードを生成させ、VS Codeで編集してLive Serverでプレビューすること。最低2回はAIにカスタマイズを依頼して改善しなさい。

\item[\textbf{演習8-3}] \textbf{GitHub Pagesで公開する}\\
演習8-2で作成したWebページをGitHub Pagesで公開しなさい。公開URLにアクセスして表示を確認し、公開URLとスクリーンショットを提出すること。

\item[\textbf{演習8-4}] \textbf{README.md の作成}\\
自分の研究プロジェクト（または架空の研究プロジェクト）について、AIを活用してREADME.mdを作成しなさい。プロジェクト概要、インストール方法、使い方（コード例付き）、ファイル構成、ライセンスを含めること。

\item[\textbf{演習8-5}] \textbf{共同作業の実践}\\
2人以上のグループで共有リポジトリを作成し、Issueを3つ以上作成して担当者を割り当て、各メンバーがブランチを作成してPull Requestとコードレビューを行いなさい。作業の流れと学んだことをレポートにまとめること。
\end{enumerate}
